The mean test error moves from \result{cl-ols-mean} for the minimum-norm OLS solution to \result{cl-ridge-mean} for ridge regression at the selected penalty, a relative reduction of \result{cl-improve} percent. This improvement is expected. In an ill-conditioned design with many more features than training rows, the minimum-norm OLS solution is forced to use directions that contain little signal, which inflates its variance. The ridge penalty shrinks these directions and trades a controlled increase in bias for a lower-variance solution. The irreducible noise floor of \result{cl-noise} gives a reference point. The regularised mean remains above that floor, so regularisation improves prediction but does not bring the model to the information limit.

However, the central limitation is that the gate certifies provenance, not validity. It verifies that a reported value matches the artifact from which it was drawn. In this study the pre-registered penalty was selected on the test set, yet the gate passed the resulting number without objection because that number was faithfully recorded. The gate did not evaluate whether selecting on the test set was a sound design choice; it only checked that the reported value was the value produced by the artifact. Provenance is not validity, and a gate that binds numbers to artifacts cannot detect a flaw in the design that produced the artifact. \citet{luo2026xscientist} makes the adjacent point that a passing integrity check is not scientific truth. The contribution is therefore narrower than an assurance of scientific truth; it guarantees faithful reporting of artifacts.
